% =====================================================================
%  Chapter 1 — The Claim
% =====================================================================
\part{The Claim}
\label{part:claim}
\markboth{Part I}{The Claim}

\chapter{What Was Actually Said}
\label{ch:claim}

On 20 November 2026 the Riemann hypothesis --- a conjecture about the zeros of
the most studied function in mathematics, open since 1859 --- acquired, briefly
and very publicly, a claimed proof. The claimant was Levent Alpöge\idx{Alpöge, Levent}, a junior
fellow at Harvard working on arithmetic geometry at Anthropic. The proof, such
as it was, was presented in two forms: a four-page document of remarkable
candour, and a Lean development of 175{,}187 files and roughly 17.4 million
lines.

This chapter reports what that document actually says. It is worth doing
carefully, because the document is not the thing most readers have encountered.
The press coverage was a distillation of a distillation, and the two most
quotable lines in the story --- ``I believe I have\ldots discovered a proof'' and
``it is a loss of a principal guiding light'' --- are not the same kind of
statement as the twenty-item technical timeline that follows them on the next
page. Both are in the source. They carry different weight, and the weight
should be assigned honestly.

\section{The four pages}The document is titled simply \emph{Remark}. It has no abstract, no theorem
statement, no definition of the objects it names, and no proof. What it has is
a narrator, a collaborator, and a defence of the decision to publish nothing
but code.

\begin{quote}
\small
``I believe I have, working with Claude\idx{Claude (AI assistant)}, discovered a proof\textsuperscript{1}
of the Riemann Hypothesis. In fact, the method easily extends much further,
resolving the so-called Generalized Riemann Hypothesis, and I expect, with more
time, the Extended Riemann Hypothesis too. Claude and I have also formalized
the arguments without assumptions. It is natural to feel that this is an end, in
my case a loss of a principal guiding light, but what I saw during the course of
progress so amazed me that I believe it is one dawn across science of something
far greater.''
\end{quote}

Three claims are made in that paragraph, and they have different strengths. The
first --- a proof of RH --- is the headline. The second --- a proof of GRH for
Dirichlet $L$-functions --- is stronger and, as the Lean artifact's own export
list shows, is what the code actually delivers. The third --- a proof of the
extended Riemann hypothesis for number fields --- is explicitly hedged with
``I expect, with more time'', and the artifact does not claim it.

The collaboration itself is stated as fact and is not in dispute. The remark
attributes the work jointly to the author and to Claude, an AI assistant, and
this book's treatment of the artifact assumes that attribution is accurate. It
also assumes, as the author says, that the human contributed the mathematical
steering: the remark describes the author finding the toy setting in which two
ideas combined, and describes Claude as having run far ahead near the end.

\section{The four pages say they are insufficient}The second paragraph of the remark is the most important sentence in the
document, and it is easy to miss because it arrives early and in a rush.

\begin{quote}
\small
``The argument is incredibly long at the moment, and while its formalization is
finished and it is routine to produce a huge, unreadable document containing
sufficient detail, in my view of my path I have just taken a step.''
\end{quote}

Read the structure of that sentence carefully. The formalisation is
\emph{finished}. A human-readable write-up of the argument \emph{is routine} and
could be produced. What is missing is not the proof but the author's
\emph{understanding} of it. And the author says so without hedging.

Page 2 makes the same point in a different register, in the middle of a
paragraph about how to publish:

\begin{quote}
\small
``This puts us in the rather extraordinary situation this age has made possible,
in which I am presenting the world with a complete argument before having
mastered a massive neighbourhood around its mathematics. It being obviously
inappropriate to attempt to post some massive unreadable travesty of `slop', I
decided to only present Lean.''
\end{quote}

Three separate concessions are packed into those sentences, and the whole
assessment of this claim turns on taking them at face value rather than
rounding them up.

\begin{enumerate}
\item The author has \emph{not mastered} the mathematics surrounding his
argument.
\item The document is explicitly a \emph{``slop'' summary} --- the author's word
for it --- of parts of the reasoning that he does not personally understand.
\item The mathematical statement of the result therefore consists of a Lean
bundle and nothing else. There is no prose description of the proof for a
referee to check, because the author says a prose description cannot yet be
written.
\end{enumerate}

There is also a footnote that fixes the scale of the artifact: ``I of course
amused myself, and hit 250,000+ pages.'' The remark is describing the rendered
Lean output, not the source.

\section{The timeline, and what it concedes}The remark's second half is a twenty-item timeline, dated, which the author
introduces with the note that ``a trusted mathematician advised me to just get to
the point and have some sort of timeline'' and that it is ``a rough timeline
Claude provided after a thorough search.'' Two entries in that timeline are
load-bearing for any assessment of the claim, and this book returns to them in
detail in Chapter~5.

\begin{seamnote}{The two entries that matter most}
\begin{itemize}
\item \textbf{14. 11/15/26:} precise computation of the ``cost'' of each lift,
showing that one would stop at 2/3 a priori, attempt to get a main term and
transfer constant to a function field or kill via twisting, need for
`division'.

\item \textbf{17. 11/15/26:} for the classes in $H^1_o(\QQ_\ell/\ZZ_\ell(-e))$
that we couldn't kill by twisting, fight for a significant weakening of
Vandiver.
\end{itemize}
\end{seamnote}

Read the dates: entry 14 and entry 17 are both 15 November. Entry 20, which
reads simply ``11/16/26: GRH'', is the following day. So the author's own record
of the work says that on 15 November the method had a computed ceiling at two
thirds and an unresolved obstruction requiring a weakening of a conjecture that
is still open; and that on 16 November, the day after, GRH was established.

The remark offers no account of how either gap was closed. It does not say. And
the remark explicitly disclaims the authority of its own summary: the parts of
the timeline in which ``Claude is the expert and not me'' are, the author
writes, ``sort of paraphrasing a `slop' summary, which might make zero sense,
and I apologize, though the Lean should suffice to clarify everything. I'm just
too inexpert in the relevant stuff at the moment.''

That last clause --- \emph{``the Lean should suffice to clarify everything''} ---
is the claim this book takes most seriously, and the one it spends the most
effort testing.

\section{The interesting sentence}One line in the remark deserves more attention than it has received, because it
is the clearest statement of the mathematical idea in the document:

\begin{quote}
\small
``In Stepanov's argument\idx{Stepanov's method} resolving Riemann over function fields one uses a
divisor involving multiplicity. This argument is a descendant of Stepanov's, but
multiplicity is hard to find\idx{multiplicity} over the integers, corresponding to something like
differentiating whole numbers.''
\end{quote}

This is a sharp diagnosis of the central difficulty, stated by someone who
understands it. The function-field proof of the Riemann hypothesis for curves
over finite fields --- Stepanov's method, refined by Bombieri --- works because
a curve over a finite field has \emph{two} independent directions to move in:
one geometric and one arithmetic. A point of multiplicity $m$ in a divisor
therefore carries $m$-dimensional information, and the method converts that
information into a bound on the number of points, which becomes a bound on the
zeros of the $L$-function.

Over the integers there is only one direction. $\operatorname{Spec}\ZZ \times
\operatorname{Spec}\ZZ = \operatorname{Spec}\ZZ$: the product is
one-dimensional. Multiplicity, in the sense Stepanov needs it, has no analogue
there. The remark's phrase ``corresponding to something like differentiating
whole numbers'' is the author reaching for an arithmetic substitute, and the
entire proposed method is an attempt to construct one.

Chapter~4 examines this barrier in detail, because it is the single most
important structural fact about the problem, and because the way the submission
answers it --- by chaining theta lifts, each of which is supposed to buy one
unit of the missing multiplicity --- is the most interesting mathematical idea
in the story.

\section{The publication decision, and its consequences}The remark describes a difficult decision. There was a timing constraint
(``the precise timing is horrible''), there was an announcement protocol to
respect, and there was a pressing practical problem: the only written form of
the proof that existed was the Lean bundle.

\begin{quote}
\small
``It being obviously inappropriate\idx{statement review} to attempt to post some massive unreadable
travesty of `slop', I decided to only present Lean.''
\end{quote}

This decision is the origin of nearly everything that makes the claim hard to
evaluate. A conventional mathematical claim comes with a statement, and the
community's first move is to check whether the statement is the interesting one.
Here the statement is a file, and the interesting statement is buried inside a
dependency graph of 175 thousand files.

It also means the claim cannot be assessed by reading. It can only be assessed
by compiling and by auditing --- which is what Chapters~6 through~9 do, and
which is why the audit in this book is technical rather than rhetorical.

\begin{keytakeaways}

\item The document under examination is a four-page narrative plus a Lean
bundle. It contains no proof, no theorem statement in prose, and no definition
of the objects it names.
\item The author states plainly that he has not mastered the surrounding
mathematics, that his own summary is ``slop'', and that the formalisation
``should suffice to clarify everything''.
\item The author's timeline records an a-priori ceiling at two thirds and an
obstruction requiring a weakening of Vandiver --- both on 15 November --- with
GRH recorded the following day and no account of how either was closed.
\item The remark's central mathematical idea --- finding an arithmetic
substitute for Stepanov's multiplicity over $\ZZ$ --- is a real and sharp
observation about where the difficulty lies.
\end{keytakeaways}

\section*{A note on how to read the rest of this book}Chapters~1 and~2 are the news and the artifact. Chapters~3 to~5 are the
mathematics, written so that a reader with undergraduate analysis can follow
them and a specialist will not find them condescending. Chapters~6 to~9 are the
audit, and are the part of the book with original work in it: the dependency
chain was traced by hand from the source tree, the scan was run over the full
extraction, and the load-bearing inputs were identified by reading them.
Chapters~10 to~12 place the claim against the literature and the prior work.
Chapters~13 to~15 give the finding, the protocol, and the consequences.

If you want only the finding, read Chapter~13 and then Chapter~14.
